\section{Conclusão}

O ruído das janelas analisadas não pode ser tratado como branco: autocorrelação
e PSD indicam correlação temporal relevante. O template do pulso concentra
energia em baixas frequências, região em que o ruído também é significativo.

Nessas condições, o filtro casado foi o elemento central para detecção e
localização do instante \(n_0\) (mediana do erro da ordem de
\(\valNZeroErrMed\) amostras). O limiar escolhido para FA aproximada de \(1\%\)
controla bem os falsos positivos, ao custo de muitos falsos negativos e F1
moderado. O filtro de Wiener, na formulação adotada, pouco alterou a
reconstrução e a detecção em relação ao sinal bruto ou ao casado puro.

A estimação de amplitude beneficia-se do BLUE que incorpora a covariância do
ruído. A reconstrução \(\hat{A}\,s[n-\hat{n}_0]\) é claramente superior ao uso
direto de \(y[n]\) ou da saída de Wiener como estimativa da forma do pulso.

Em síntese: para detectar, o filtro casado; para estimar amplitude e
reconstruir a forma, o BLUE colorido alinhado pelo casado. Os resultados no
teste acompanham os do treino, sem retreino de limiar.

\section*{Entrega}
\addcontentsline{toc}{section}{Entrega}

São entregues dois arquivos: este relatório em PDF e o caderno Jupyter com a
análise completa. O PDF reúne a teoria dos filtros, os gráficos, as estimativas
e a interpretação dos resultados. O caderno contém o código utilizado para
carregar os dados, projetar os filtros, ajustar o limiar no treino e avaliar o
desempenho no teste.
